\scanpage{150}
partie de $***$. Mais si \deleted{\unsure{$u\geq0$}}\inserted{$\unsure{1|-1}$ n'est pas valeur propre,} on a aussi, \deleted{pour tout} \inserted{\unsure{pour la dimension infinie|particulièrement en dimension infinie}} $\rho_i\deleted{\unsure{\geq}}=\frac{\det(1+\unsure{\lambda|u})}{1+\lambda_i}$ \deleted{\unsure{$\geq$|$\leq$}} d'où

\[
\boxed{\|R(u)\|_\infty\geq|\det(1+u)|}\quad\text{(\unsure{op normal|si dim infinie})}
\]

(car il existe des $\lambda_i$ arbitrairement petits, donc des $\rho_i$ arbitrairement \deleted{distants \unsure{par}}\inserted{voisins} de $\det(1+u)$ !

On \deleted{\unsure{en}} prouve donc que l'inégalité $***$ a) ne peut s'améliorer, car elle devient égalité pour $u\geq0$ !
% archive sheet number: 151
